\documentclass[11pt,a4paper]{report}

% ------------------------------------------------------------------
% Informe: IEEE-CIS Fraud Detection — ingeniería de datos y modelos
% Compilación: tectonic main.tex  (o xelatex main.tex ×2)
% ------------------------------------------------------------------
\usepackage[spanish,es-tabla]{babel}
\usepackage[margin=2.5cm]{geometry}
\usepackage{amsmath,amssymb}
\usepackage{graphicx}
\usepackage{booktabs}
\usepackage{multirow}
\usepackage{subcaption}
\usepackage{xcolor}
\usepackage{listings}
\usepackage{microtype}
\usepackage{csquotes}
\usepackage[colorlinks=true,linkcolor=blue!60!black,citecolor=green!40!black,urlcolor=blue!60!black]{hyperref}
\usepackage[nameinlink,capitalize]{cleveref}

% --- listings estilo Python ---
\definecolor{codebg}{RGB}{248,248,248}
\definecolor{codekey}{RGB}{0,92,197}
\definecolor{codestr}{RGB}{163,21,21}
\lstdefinestyle{py}{
  language=Python,
  basicstyle=\ttfamily\footnotesize,
  backgroundcolor=\color{codebg},
  keywordstyle=\color{codekey}\bfseries,
  stringstyle=\color{codestr},
  commentstyle=\color{gray}\itshape,
  showstringspaces=false,
  breaklines=true,
  frame=single,
  rulecolor=\color{gray!40},
  columns=fullflexible,
  literate={á}{{\'a}}1 {é}{{\'e}}1 {í}{{\'i}}1 {ó}{{\'o}}1 {ú}{{\'u}}1 {ñ}{{\~n}}1
}
\lstset{style=py}

% --- accesos rápidos ---
\newcommand{\figdir}{figures/}
\newcommand{\tabdir}{tables/}

\title{\textbf{Ingeniería de datos y modelos de clasificación para\\ detección de fraude en transacciones}\\[6pt]
\large Estudio comparativo sobre la competencia IEEE-CIS Fraud Detection\\
\large Proyecto 1 — Planificación}
\author{Edgar Chambilla}
\date{Septiembre de 2026}

\begin{document}
\maketitle

\begin{abstract}
\noindent
Este informe documenta la reproducción local y el análisis comparativo de las mejores soluciones
públicas de la competencia \emph{IEEE-CIS Fraud Detection} (Kaggle, 2019), con el objetivo de
identificar qué técnicas de tratamiento de datos y qué modelos de clasificación explican el
rendimiento observado. Se seleccionaron los 12 notebooks públicos con mayor reputación, se
descargaron sus fuentes y se ejecutaron localmente 4 de ellos (los que producen modelos
entrenables) sobre el conjunto completo de datos, bajo validación temporal que replica la
separación train/test de la competencia. El resultado central es una \emph{ablación} controlada:
con exactamente el mismo modelo (XGBoost, profundidad 12), la construcción de un pseudo-identificador
de cliente (\emph{magic UID}) y sus $\sim$40 agregaciones por grupo elevan el AUC local de
0{,}9375 a 0{,}9479 ($+0{,}0104$), mientras que los cambios de algoritmo o de hardware apenas
movilizan el AUC. Se concluye que, en problemas con distribución temporal no estacionaria,
la jerarquía de impacto es: (1)~validación honesta, (2)~ingeniería de features orientada a
entidades, (3)~estabilización temporal de features, y (4)~recién entonces, hiperparámetros y
arquitectura del modelo. Todos los experimentos son reproducibles con los scripts incluidos.
\end{abstract}

\tableofcontents
\listoffigures
\listoftables

\chapter{Introducción}
\label{cap:intro}

\section{Motivación}

El fraude en pagos \emph{card-not-present} (CNP) genera pérdidas anuales de miles de millones de
dólares y plantea un problema de clasificación binaria particular: el conjunto de datos está
fuertemente desbalanceado (alrededor del 3{,}5\,\% de transacciones fraudulentas), y las features
son en su mayoría anónimas o están ofuscadas por razones comerciales. El futuro tampoco se parece al
pasado: los patrones de ataque cambian, por lo que un modelo que
sobreajusta el período de entrenamiento degrada rápidamente su desempeño en producción.

La competencia \emph{IEEE-CIS Fraud Detection}~\cite{ieee-cis} (Kaggle, 2019), organizada por la
IEEE Computational Intelligence Society junto a Vesta Corporation, formalizó este problema sobre
transacciones reales de comercio electrónico. Contó con 6\,355 equipos participantes y consolidó
un cuerpo público de soluciones, entre las que destaca la técnica del
\emph{magic UID} de Chris Deotte y Konstantin Yakovlev (equipo \emph{FraudSquad}, 2.\textsuperscript{o}
en el tablero público y 1.\textsuperscript{o} en el privado).

\section{Objetivos}

\begin{enumerate}
  \item \textbf{Recolectar} las mejores soluciones públicas de la competencia (notebooks con
        mayor reputación y write-ups de equipos ganadores).
  \item \textbf{Reproducir localmente} los notebooks que producen modelos entrenables, mediendo
        AUC y costo computacional sobre hardware propio (16 núcleos CPU, NVIDIA RTX 2060 de 6\,GB).
  \item \textbf{Aislar} mediante ablación controlada qué componente explica la ganancia de
        rendimiento: el algoritmo, el hardware o la ingeniería de variables.
  \item \textbf{Documentar} las técnicas de tratamiento de datos y de modelado, dejando un
        pipeline reproducible como base para una propuesta propia.
\end{enumerate}

\section{Metodología general}

El estudio siguió cuatro fases:

\begin{enumerate}
  \item \textbf{Selección de código}: se identificaron los 12 notebooks públicos con mayor
        reputación (2\,986 a 169 votos) mediante el índice curado \emph{Must-see Kernels}
        \cite{mustsee} y el API público de Kaggle. Cada fuente se descargó en formato
        \texttt{.ipynb} y se exportó a Python legible.
  \item \textbf{Triaje}: se clasificó cada notebook en \emph{ejecutable} (produce un modelo
        entrenable) o \emph{analítico} (EDA pura, técnicas de re-muestreo, infraestructura GPU
        no compatible con 6\,GB de VRAM). Cuatro notebooks resultaron ejecutables.
  \item \textbf{Reproducción local}: los 4 ejecutables se adaptaron a las versiones modernas de
        las librerías (XGBoost~$\ge$~2, LightGBM~$\ge$~4, pandas~3) y se ejecutaron con el
        conjunto completo de entrenamiento, bajo un protocolo de validación temporal uniforme
        (\cref{sec:validacion}).
  \item \textbf{Análisis comparativo}: se comparó AUC y tiempo de cómputo, y se realizó una
        ablación \emph{con/sin magic UID} sobre el mismo modelo para atribuir la ganancia.
\end{enumerate}

\section{Estructura del informe}

\cref{cap:datos} describe el conjunto de datos; \cref{cap:tratamiento} detalla las técnicas de
tratamiento y de ingeniería de variables; \cref{cap:modelos} presenta los modelos comparados;
\cref{cap:resultados} muestra los resultados locales con sus gráficos; \cref{cap:discusion}
discute ventajas, límites metodológicos y lecciones; \cref{cap:conclusiones} concluye y propone
el desarrollo futuro. \cref{cap:apendice} documenta la reproducibilidad completa.

\chapter{Conjunto de datos}
\label{cap:datos}

\section{Origen y estructura}

Los datos provienen de transacciones reales de comercio electrónico provistas por Vesta
Corporation. Se componen de dos tablas unidas por \texttt{TransactionID}
(\cref{tab:datos}): una tabla de \emph{transacción} (590\,540 filas $\times$ 394 columnas en
entrenamiento) y una tabla de \emph{identidad} (144\,233 filas $\times$ 40 columnas), disponible
solo para el $\sim$24\,\% de las transacciones.

\begin{table}[ht]
  \centering
  \caption{Resumen de los archivos del conjunto de datos.}
  \label{tab:datos}
  \input{\tabdir tabla_datos}
\end{table}

Las variables se agrupan en familias semánticas:

\begin{itemize}
  \item \textbf{Montos y producto}: \texttt{TransactionAmt}, \texttt{ProductCD} (W, C, R, H, S).
  \item \textbf{Tarjeta}: \texttt{card1}--\texttt{card6} (emisor, tipo, categoría).
  \item \textbf{Geografía}: \texttt{addr1}--\texttt{addr2}, \texttt{dist1}--\texttt{dist2}.
  \item \textbf{Correo}: \texttt{P\_emaildomain}, \texttt{R\_emaildomain}.
  \item \textbf{Conteos} \texttt{C1}--\texttt{C14}: cuántos elementos asociados (tarjetas,
        direcciones, correos) aparecen en la ventana de observación.
  \item \textbf{Deltas de tiempo} \texttt{D1}--\texttt{D15}: días transcurridos desde un evento
        previo (p.\,ej.\ primera transacción, cambio de tarjeta).
  \item \textbf{Coincidencias} \texttt{M1}--\texttt{M9}: variables booleanas de matching.
  \item \textbf{\texttt{V1}--\texttt{V339}}: variables propietarias de Vesta (rangos, conteos,
        hashes), con bloques fuertemente correlacionados.
  \item \textbf{Identidad} \texttt{id\_01}--\texttt{id\_38}, \texttt{DeviceType},
        \texttt{DeviceInfo}: huella de dispositivo y de red.
\end{itemize}

\section{Desbalance del problema}

De las 590\,540 transacciones de entrenamiento, 20\,663 son fraudulentas: un 3{,}50\,\%
(razón de desbalance 1:27, \cref{fig:target}). Esta proporción hace que la exactitud
(\emph{accuracy}) sea una métrica inútil ---un clasificador trivial que predice ``legítima''
alcanza 96{,}5\,\% de exactitud con utilidad nula---, por lo que la competencia adoptó el
área bajo la curva ROC (AUC) como métrica oficial, que ordena transacciones en lugar de
etiquetarlas con umbral fijo.

\begin{figure}[ht]
  \centering
  \includegraphics[width=0.62\textwidth]{\figdir fig_distribucion_target.pdf}
  \caption{Distribución de la variable objetivo. El desbalance 1:27 exige métricas de orden
           (AUC) y hace innecesario el re-muestreo agresivo para modelos de árboles.}
  \label{fig:target}
\end{figure}

\section{Calidad de los datos}

Dos problemas dominan la calidad de los datos. Primero, los \textbf{valores faltantes}: la
familia \texttt{V} promedia un 43\,\% de faltantes y la tabla de identidad, al fusionarse por
completitud izquierda, introduce un 76\,\% de faltantes estructurales en sus columnas
(\cref{fig:nulos}). Segundo, la \textbf{cola pesada} del monto (\cref{fig:amt}), que en escala
lineal concentra el 99\,\% de la masa por debajo de \$300 mientras la cola llega a \$31\,937:
los montos extremos son justamente sospechosos, por lo que no deben recortarse sino
re-escalarse (logaritmo) o dejarse en bruto para los árboles.

\begin{figure}[ht]
  \centering
  \includegraphics[width=0.85\textwidth]{\figdir fig_nulos_por_grupo.pdf}
  \caption{Porcentaje promedio de valores faltantes por familia de variables. La ausencia de la
           tabla de identidad es estructural (transacciones sin huella de dispositivo) y, como se
           verá, codificar la ausencia como categoría propia (-999) es preferible a imputar.}
  \label{fig:nulos}
\end{figure}

\begin{figure}[ht]
  \centering
  \includegraphics[width=0.85\textwidth]{\figdir fig_amt_distribucion.pdf}
  \caption{Distribución del monto por clase en escala logarítmica. Las transacciones fraudulentas
           se concentran en montos bajos-repetitivos, con una cola superior más densa que la de
           las legítimas.}
  \label{fig:amt}
\end{figure}

\section{La estructura temporal: la dificultad central}

\texttt{TransactionDT} es un desplazamiento relativo en segundos (693\,999--1\,581\,113\,131)
que cubre 182 días de entrenamiento. La partición oficial de la competencia es
\textbf{temporal}: el conjunto de prueba corresponde al mes siguiente al final del
entrenamiento. Esto tiene tres consecuencias metodológicas:

\begin{enumerate}
  \item Una validación aleatoria (KFold clásico) \emph{sobreestima} el desempeño: filas
        ``del futuro'' aparecen en el entrenamiento de cada fold.
  \item Cualquier variable cuya distribución derive con el tiempo produce una relación espuria
        con el target (\cref{fig:fraude-temporal} y \cref{sec:dcols}).
  \item La validación honesta debe simular el futuro: \emph{holdout} por corte temporal,
        \texttt{TimeSeriesSplit}, o \texttt{GroupKFold} agrupando por mes calendario.
\end{enumerate}

\begin{figure}[ht]
  \centering
  \includegraphics[width=0.95\textwidth]{\figdir fig_fraude_temporal.pdf}
  \caption{Tasa diaria de fraude durante los 182 días de entrenamiento. Los picos agudos
           corresponden a campañas de ataque acotadas en el tiempo: el patrón no es estacionario,
           y por eso la selección de variables por \emph{consistencia temporal}
           (\cref{sec:timeconsistency}) resulta decisiva.}
  \label{fig:fraude-temporal}
\end{figure}

Como referencia de mercado, \cref{fig:leaderboard} muestra la distribución de puntajes del
tablero público (6\,355 equipos): la masa de competidores se concentra bajo 0{,}95 y el top-5
supera 0{,}966, lo que ilustra el diferencial que aportan las técnicas que este informe estudia.

\begin{figure}[ht]
  \centering
  \includegraphics[width=0.9\textwidth]{\figdir fig_leaderboard.pdf}
  \caption{Histograma del tablero público de la competencia. En rojo, los cinco primeros
           equipos: AlKo (0{,}9681), FraudSquad (0{,}9677), Young for you (0{,}9676),
           2 uncles and 3 puppies (0{,}9672) y Mr Lonely (0{,}9663).}
  \label{fig:leaderboard}
\end{figure}

\chapter{Tratamiento de datos e ingeniería de variables}
\label{cap:tratamiento}

Este capítulo documenta las técnicas de preparación de datos extraídas del código real de los
cuatro notebooks ejecutados, ordenadas de lo básico a lo que diferencia a la solución ganadora.

\section{Integración y tipos de datos}

Las dos tablas se integran por \texttt{TransactionID} con unión izquierda (la identidad es
opcional). Dado que \texttt{train\_transaction.csv} ocupa 2{,}2\,GB en memoria por defecto, se
aplican dos optimizaciones de tipo:

\begin{itemize}
  \item numéricos a \texttt{float32} / \texttt{int16} donde el rango lo permite (reduce la
        memoria a la mitad, habilitando GPU);
  \item categóricas a \texttt{category} (diccionario de enteros) o \texttt{int16} vía
        \texttt{factorize}.
\end{itemize}

\section{Valores faltantes: codificar en vez de imputar}

Los notebooks de referencia convierten el faltante en \emph{información}:

\begin{lstlisting}
X_train = X_train.fillna(-999)          # baseline (inversion, xhlulu)
clf = xgb.XGBClassifier(..., missing=-999)
\end{lstlisting}

El árbol de decisión aprende que $-999$ es una ``categoría desconocida'' y puede ramificar
sobre ella; imputar con la media, en cambio, fabrica valores plausibles y falsos. En la
solución de Deotte el esquema es más fino: los numéricos se desplazan para hacerlos positivos
y el faltante se codifica como $-1$:
\begin{lstlisting}
mn = min(X_train[f].min(), X_test[f].min())
X_train[f] -= mn;  X_train[f].fillna(-1, inplace=True)
\end{lstlisting}

\section{Codificación de variables categóricas}

Tres variantes aparecen en el código estudiado:

\begin{enumerate}
  \item \textbf{Etiqueta entera} (\texttt{LabelEncoder} / \texttt{factorize}): suficiente para
        árboles, que tratan los códigos como categorías ordinales arbitrarias. Se ajusta sobre
        train y test juntos para compartir el diccionario.
  \item \textbf{Codificación de frecuencia} (\emph{count encoding}, nroman): reemplaza cada
        categoría por su frecuencia de aparición.
        \begin{lstlisting}
train['card1_count'] = train['card1'].map(train['card1'].value_counts(dropna=False))
        \end{lstlisting}
        Es una agregación ``de bolsillo'': una tarjeta con miles de transacciones es una tarjeta
        comprometida, y esta señal queda expuesta como número.
  \item \textbf{Interacciones como cadena} (nroman, Deotte): concatenar dos categóricas crea una
        entidad más fina.
        \begin{lstlisting}
train['card1__card5'] = train['card1'].astype(str) + '_' + train['card5'].astype(str)
train['uid'] = train['card1_addr1'].astype(str) + '_' + (day - D1).astype(str)
        \end{lstlisting}
\end{enumerate}

Variables derivadas simples pero efectivas: día de la semana y hora del día a partir de
\texttt{TransactionDT}; parte decimal del monto ($(A - \lfloor A\rfloor)\times 1000$, patrón de
centavos de bots); \texttt{log1p} del monto (kyakovlev).

\section{Normalización de columnas D (estabilización temporal)}
\label{sec:dcols}

Las columnas \texttt{D1}--\texttt{D15} significan ``días transcurridos desde un evento previo''
medidos \emph{en el momento de la transacción}: por construcción derivan con el tiempo absoluto
y el modelo confunde ``mucho tiempo transcurrido'' con ``fraude''. La solución de Deotte
resta el reloj absoluto:

\begin{lstlisting}
for i in [4,5,6,8,10,11,12,13,14,15]:
    X_train['D'+str(i)] -= X_train.TransactionDT / (24*60*60)
    X_test ['D'+str(i)] -= X_test.TransactionDT  / (24*60*60)
\end{lstlisting}

\cref{fig:d15} muestra el efecto sobre \texttt{D15}: de una nube con pendiente (train y test
viven en rangos distintos) a una banda horizontal estacionaria, directamente transferible al
futuro. La misma idea es la que motiva la \emph{validación adversarial} presente en el corpus
estudiado: entrenar un clasificador train-vs-test y normalizar las variables que lo separan.

\begin{figure}[ht]
  \centering
  \includegraphics[width=0.98\textwidth]{\figdir fig_d15_normalizacion.pdf}
  \caption{Normalización de \texttt{D15}. Izquierda: la variable cruda deriva con el tiempo.
           Derecha: tras restar el reloj absoluto queda una banda estable, comparable entre
           entrenamiento y futuro.}
  \label{fig:d15}
\end{figure}

\section{Selección por consistencia temporal}
\label{sec:timeconsistency}

La solución de Deotte elimina explícitamente las variables cuyo histograma difiere entre el
inicio y el fin del período de entrenamiento (p.\,ej.\ \texttt{C3}, \texttt{M5}, \texttt{id\_08},
\texttt{id\_33}, \texttt{card4}). El criterio es estadístico y barato: si una variable ya no es
estable dentro del entrenamiento, su relación con el fraude tampoco lo será en el mes de prueba.
De las 339 columnas \texttt{V} se conservan $\sim$120, elegidas por bloques de correlación en un
EDA previo del mismo autor.

\section{El pseudo-identificador de cliente (\emph{magic UID})}
\label{sec:uid}

El dato que más ayudaría ---el identificador del cliente--- no existe en el conjunto. La
solución ganadora lo reconstruye a partir de tres columnas:
\begin{equation}
  \texttt{uid} \;=\; \underbrace{\texttt{card1\_addr1}}_{\text{tarjeta + dirección}}
  \;\Vert\;
  \underbrace{\big\lfloor \tfrac{\texttt{TransactionDT}}{86400} - \texttt{D1} \big\rfloor}_{\text{día estimado de alta de la cuenta}}
  \label{eq:uid}
\end{equation}
donde \texttt{D1} (días desde la primera transacción) actúa como huella de antigüedad de la
cuenta. La combinación sobrevive al cambio de dispositivo y de correo. \cref{fig:uid} resume el
esquema completo: a partir del UID se generan tres familias de agregaciones y \emph{luego el UID
se elimina}; nunca entra al modelo como variable.

\begin{figure}[ht]
  \centering
  \includegraphics[width=0.98\textwidth]{\figdir fig_esquema_uid.pdf}
  \caption{Esquema de la técnica del \emph{magic UID}: reconstrucción del cliente a partir de
           \eqref{eq:uid} y generación de agregaciones por grupo. El identificador se usa solo
           como llave de agrupación y se descarta antes del entrenamiento.}
  \label{fig:uid}
\end{figure}

Las tres familias de agregaciones, tal como aparecen en el código original:

\begin{lstlisting}
# frecuencia del grupo: ¿cuántas transacciones acumula este cliente?
encode_FE(X_train, X_test, ['uid'])

# media y desviacion por grupo: ¿es anomalo este monto/delta PARA ESTE cliente?
encode_AG(['TransactionAmt','D4','D9','D10','D15'], ['uid'], ['mean','std'])
encode_AG(['C'+str(x) for x in range(1,15) if x!=3], ['uid'], ['mean'])
encode_AG(['M'+str(x) for x in range(1,10)], ['uid'], ['mean'])

# cardinalidad por grupo: ¿cuantos valores distintos usa el cliente?
encode_AG2(['P_emaildomain','dist1','id_02','cents'], ['uid'])
encode_AG2(['C13','V314','V127','V136','V309','V307','V320'], ['uid'])

# interaccion ad-hoc: direccion de envio de otro "mundo" temporal
X_train['outsider15'] = (np.abs(X_train.D1 - X_train.D15) > 3).astype('int8')
\end{lstlisting}

La intuición estadística es la reducción de varianza intra-grupo: un árbol que solo ve la
variable $X$ comete errores de contexto, mientras que la media de $X$ dentro del grupo del
cliente contextualiza cada observación y vuelve trivial la separación de los casos anómalos.
\cref{sec:ablacion} cuantifica este efecto en el experimento de ablación: $+0{,}0104$ de AUC
con el modelo idéntico.

\section{Protocolo de validación}
\label{sec:validacion}

Dada la partición temporal de la competencia, se adoptaron tres esquemas según el propósito:

\begin{description}
  \item[Holdout temporal 75/25] entrenar con los primeros $\sim$136 días y validar con los
        últimos $\sim$46. Rápido, para iterar (baseline, XGB\_95/96).
  \item[\texttt{TimeSeriesSplit} $\times$3] ventanas crecientes con validación siempre
        ``hacia adelante'' (nroman).
  \item[\texttt{GroupKFold} mensual $\times$3] cada fold reserva \emph{meses completos} como
        validación, agrupando por mes calendario reconstruido desde
        \texttt{START\_DATE 2017-11-30} (XGB\_95/96). Es el esquema más fiel al operativo real
        y el que se reporta como AUC local principal.
\end{description}

\section{Limitación conocida del protocolo}

Las agregaciones por UID se computan sobre \emph{todo} el período de entrenamiento, incluido el
mes reservado como validación (aunque \emph{sin} usar la etiqueta del mes reservado para
construir el modelo: el leak es de estadísticas de grupo, no de clases). Esto reproduce
fielmente el protocolo original de la competencia y explica parte del nivel absoluto de AUC
local; una variante estrictamente causal (agregar solo transacciones anteriores a cada
\texttt{TransactionDT}) queda propuesta como trabajo futuro en \cref{sec:futuro}.

\chapter{Modelos comparados}
\label{cap:modelos}

Los cuatro notebooks ejecutables comparten la misma familia algorítmica ---\emph{gradient
boosted trees}--- y difieren en el tratamiento de datos, el hardware y el protocolo de
validación. Esta homogeneidad es intencional: hace atribuible cualquier diferencia de AUC a la
ingeniería de variables. \cref{tab:modelos} resume las configuraciones.

\begin{table}[ht]
  \centering
  \caption{Configuraciones comparadas. ``d9/d12'' denota profundidad máxima del árbol.}
  \label{tab:modelos}
  \small
  \begin{tabular}{lllll}
    \toprule
    & Baseline CPU & Baseline GPU & LightGBM (nroman) & XGB d12 (cdeotte) \\
    \midrule
    Librería        & XGBoost 3.4  & XGBoost 3.4  & LightGBM 4.x & XGBoost 3.4 \\
    Dispositivo     & CPU (16 hilos) & RTX 2060    & CPU (16 hilos) & RTX 2060 \\
    Profundidad     & 9            & 9            & 31 hojas     & 12 \\
    Árboles         & 500          & 500          & $\le$5\,000 (ES) & 2\,000--5\,000 (ES) \\
    Tasa de aprend. & 0.05         & 0.05         & 0.0069       & 0.02 \\
    Features        & 432 crudas   & 432 crudas   & 298 (RFE + FE) & $\sim$200 (V + UID) \\
    Categóricas     & etiqueta     & etiqueta     & etiqueta     & etiqueta + frecuencia \\
    Validación      & holdout 75/25 & holdout 75/25 & TSSplit $\times$3 & GroupKFold mes $\times$3 \\
    \bottomrule
  \end{tabular}
\end{table}

\section{Baseline XGBoost (inversion, xhlulu)}

El punto de referencia: unión de tablas, relleno de faltantes con $-999$, etiquetado entero de
categóricas y \texttt{XGBClassifier} de profundidad 9 con 500 árboles. La variante GPU delega el
histograma y la construcción de árboles a la RTX 2060. En XGBoost moderno la invocación es:
\begin{lstlisting}
xgb.XGBClassifier(n_estimators=500, max_depth=9, learning_rate=0.05,
                  subsample=0.9, colsample_bytree=0.9, missing=-999,
                  tree_method='hist', device='cuda')   # ex 'gpu_hist'
\end{lstlisting}

\section{LightGBM con selección de variables (nroman)}

Parte de 120 variables ``útiles'' obtenidas por eliminación recursiva de características (RFE)
en un notebook hermano, añade codificación de frecuencia, variables de calendario,
interacciones y el decimal del monto, y entrena LightGBM con hiperparámetros de optimización
bayesiana (\cref{tab:modelos}), validación \texttt{TimeSeriesSplit} y parada temprana. El modelo
final re-entrena con todo el período y el número óptimo de iteraciones.

\section{XGBoost con ingeniería de entidades (cdeotte XGB\_95 / XGB\_96)}

Un único diseño de modelo (profundidad 12, 2\,000--5\,000 árboles, tasa 0.02, muestreo
0.8/0.4) evaluado en dos niveles de tratamiento de datos:

\begin{description}
  \item[XGB\_95 (sin UID):] selección de $\sim$120 columnas \texttt{V} por bloques de
        correlación, normalización de columnas \texttt{D}, codificación de frecuencia de
        tarjeta/dirección/correo, combinaciones \texttt{card1+addr1(+email)}, agregaciones
        \emph{por tarjeta} de monto y \texttt{D9}/\texttt{D11}, y descarte por consistencia
        temporal.
  \item[XGB\_96 (con magic UID):] todo lo anterior \emph{más} el pseudo-cliente de
        \cref{sec:uid} y sus $\sim$40 agregaciones. El modelo es idéntico al de XGB\_95:
        la ablación de \cref{sec:ablacion} aísla el aporte del UID.
\end{description}

El notebook original cierra con un post-procesado que suaviza las probabilidades promediando
dentro de cada UID (transacciones del mismo cliente deberían recibir scores similares); ese
paso requiere listas externas de limpieza de UIDs y queda fuera del alcance local, documentado
como trabajo futuro.

\section{Costo computacional}

La RTX 2060 (6\,GB) acelera el entrenamiento 4.4$\times$ respecto a 16 núcleos CPU en el
baseline, con AUC indistinguible: la GPU es una palanca de \emph{productividad} (más
iteraciones por hora), no de precisión. El límite de memoria también define fronteras
prácticas: el flujo \emph{RAPIDS cuDF} del mismo autor (todo el preprocesamiento en GPU,
15$\times$ más rápido) exige más de 6\,GB de VRAM y no fue ejecutable en este hardware;
queda documentado en el corpus descartado.

\chapter{Resultados}
\label{cap:resultados}

\section{Comparación principal}

\cref{tab:ranking} presenta los resultados de la reproducción local. Todos los AUC provienen de
validación temporal ejecutada en esta máquina sobre el conjunto completo de entrenamiento; no se
realizó ninguna submisión a Kaggle (el estudio es estrictamente local).

\begin{table}[ht]
  \centering
  \caption{Ranking local: AUC y costo computacional por método. El AUC principal de cada método
           se reporta en su esquema de validación propio (columna Validación).}
  \label{tab:ranking}
  \resizebox{\textwidth}{!}{\input{\tabdir tabla_ranking}}
\end{table}

\begin{figure}[ht]
  \centering
  \includegraphics[width=0.9\textwidth]{\figdir fig_auc_comparacion.pdf}
  \caption{AUC local por método. El salto relevante ocurre al pasar de features crudas a la
           ingeniería de entidades (rojo): el mismo algoritmo gana $+0{,}018$ con la preparación
           de datos de Deotte y $+0{,}010$ adicionales con el magic UID.}
  \label{fig:auc-comp}
\end{figure}

\begin{figure}[ht]
  \centering
  \includegraphics[width=0.88\textwidth]{\figdir fig_auc_vs_tiempo.pdf}
  \caption{Precisión vs.\ costo computacional. La GPU reduce el tiempo del baseline de 112\,s a
           25\,s (4.4$\times$) sin alterar el AUC; la mayor precisión exige más cómputo
           (profundidad 12 y hasta 5\,000 árboles) pero la ganancia dominante proviene de los
           datos, no del cómputo.}
  \label{fig:auc-tiempo}
\end{figure}

\section{Ablación del magic UID}
\label{sec:ablacion}

\cref{fig:ablacion} muestra el experimento clave. Con el modelo, los hiperparámetros y el
hardware idénticos, añadir el pseudo-cliente y sus agregaciones produce:

\begin{equation}
  \Delta\text{AUC}_{\text{UID}} = 0{,}9479 - 0{,}9375 = +0{,}0104
\end{equation}
medido con validación OOF \texttt{GroupKFold} mensual.

La ganancia es superior a la de cualquier otro cambio estudiado y se atribuye íntegramente a la
representación de datos: el identificador nunca ingresa al árbol, solo sus estadísticas de
grupo. En el holdout temporal 75/25 el efecto es consistente ($0{,}9433$ vs.\ $0{,}9370$).

\begin{figure}[ht]
  \centering
  \includegraphics[width=0.78\textwidth]{\figdir fig_ablacion_uid.pdf}
  \caption{Ablación con/sin magic UID sobre el mismo modelo (XGBoost d12). La diferencia entre
           barras del mismo color es el aporte puro de la ingeniería de entidades.}
  \label{fig:ablacion}
\end{figure}

\section{Qué variables dominan}

Un ajuste exploratorio rápido de XGBoost (150 árboles, sobre las familias \texttt{C/D/V} y las
columnas de tarjeta) confirma la estructura de señal esperada: dominan los conteos
\texttt{C13}, \texttt{C14}, \texttt{C1}, los deltas de tiempo normalizables (\texttt{D10},
\texttt{D15}, \texttt{D4}) y las columnas de tarjeta/dirección que componen el UID
(\cref{fig:topfeat}). Es coherente con la tesis del informe: la señal discriminante es
\emph{comportamental} (actividad por entidad y ventanas de tiempo), no transaccional aislada.

\begin{figure}[ht]
  \centering
  \includegraphics[width=0.82\textwidth]{\figdir fig_top_features.pdf}
  \caption{Importancia (gain) de las 20 variables principales en un XGBoost exploratorio.
           El peso de conteos, deltas temporales y tarjetas anticipa por qué las agregaciones
           por cliente resultan tan efectivas.}
  \label{fig:topfeat}
\end{figure}

\section{Referencias cruzadas con Kaggle}

Los puntajes públicos de los mismos métodos en la competencia (fuera de este estudio, solo como
referencia) fueron: baseline $\sim$0{,}943, nroman 0{,}9419, XGB\_95 0{,}9514 y XGB\_96
0{,}9627. El \emph{ordenamiento} local coincide con el de la competencia: más ingeniería de
entidades $\Rightarrow$ más AUC. El nivel absoluto local es menor porque el tablero evalúa el
mes posterior con un cuarto de las transacciones y los puntajes públicos históricos se obtuvieron
con submisión directa; nuestra validación local es deliberadamente más exigente.

\chapter{Discusión}
\label{cap:discusion}

\section{Ventajas y límites de cada método}

\begin{table}[ht]
  \centering
  \caption{Ventajas y limitaciones observadas en la reproducción local.}
  \label{tab:ventajas}
  \small
  \begin{tabular}{p{0.30\textwidth}p{0.32\textwidth}p{0.30\textwidth}}
    \toprule
    Método & Ventaja principal & Límite principal \\
    \midrule
    Baseline XGBoost CPU &
      Mínimo costo de desarrollo; piso de referencia inmediato ($\sim$0.92). &
      Features crudas desperdician la señal de entidad; 112\,s por ajuste frena la iteración. \\
    \addlinespace
    Baseline XGBoost GPU &
      4.4$\times$ más rápido con AUC idéntico; habilita búsqueda de hiperparámetros. &
      No aporta precisión por sí solo; exige VRAM suficiente. \\
    \addlinespace
    LightGBM + RFE + count encoding (nroman) &
      30\,\% menos features con AUC equivalente; validación temporal estricta; interpretable. &
      Ganancia de precisión marginal; la tasa de aprendizaje 0.0069 exige miles de árboles. \\
    \addlinespace
    XGBoost d12 + magic UID (cdeotte) &
      Mayor AUC local (0.9479); ganancia atribuible y replicable ($+0.0104$). &
      Requiere conocimiento del dominio (composición del UID); agregaciones costosas; protocolo
      transductivo (ver \cref{sec:limites}). \\
    \bottomrule
  \end{tabular}
\end{table}

\section{Qué importa realmente: jerarquía de impacto}

Los experimentos permiten ordenar los factores por su contribución al AUC:

\begin{enumerate}
  \item \textbf{Validación honesta} (no aporta AUC, evita ilusiones): el esquema de validación
        decide si un resultado es real. KFold aleatorio habría reportado valores notablemente
        superiores e irrepetibles en producción.
  \item \textbf{Ingeniería de entidades} ($+0{,}028$ acumulado): el pseudo-cliente y sus
        agregaciones explican la mayor parte de la distancia entre el baseline y el tope local.
        La lección es transferible a cualquier dominio con entidades implícitas (usuarios,
        dispositivos, comercios): \emph{construir la entidad, agregar su historial, descartar la
        llave}.
  \item \textbf{Estabilización temporal} (habilitadora): normalizar columnas \texttt{D} y filtrar
        por consistencia temporal no sube el AUC por sí sola, pero hace que las ganancias
        anteriores \emph{se transfieran} al futuro.
  \item \textbf{Selección de variables} ($\approx$neutral aquí): recortar de 432 a 120 variables
        mantuvo el AUC con menos costo, relevante en producción.
  \item \textbf{Algoritmo e hiperparámetros} (menor): entre XGBoost y LightGBM la diferencia
        observada es inferior al efecto de cualquier técnica de datos; la optimización bayesiana
        de nroman vale menos que una sola feature de entidad bien construida.
  \item \textbf{Hardware} ($0$ AUC, $\times$4.4 velocidad): la GPU acelera el ciclo de
        experimentación, que es donde se generan las ganancias de datos.
\end{enumerate}

\section{Limitaciones metodológicas}
\label{sec:limites}

\begin{enumerate}
  \item \textbf{Protocolo transductivo}: las agregaciones por UID se computan sobre todo el
        período de entrenamiento, incluidos los meses usados como validación (sin usar sus
        etiquetas). El AUC local absoluto es, por tanto, un límite superior respecto de un
        sistema causal estricto. La variante estricta (estadísticas solo con transacciones
        anteriores a cada transacción) es el siguiente hito propuesto.
  \item \textbf{Un solo grano de validación para los XGB d12}: GroupKFold de 3 meses (6 en el
        original) por restricciones de tiempo de GPU; las diferencias entre folds fueron
        pequeñas y consistentes en dirección.
  \item \textbf{Corpus no ejecutado}: 8 notebooks de alto valor analítico (EDA, re-muestreo,
        optimización bayesiana completa, \emph{blend} de artgor, RAPIDS de Deotte) se analizaron
        cualitativamente pero no se reprodujeron numéricamente; sus conclusiones no cuentan con
        medición local.
  \item \textbf{Métrica única}: el AUC ordena bien pero no responde por costos asimétricos de
        negocio (falso positivo bloquea un pago legítimo); un despliegue real requiere umbral y
        métrica de utilidad propios.
\end{enumerate}

\section{Amenazas a la validez}

Las adaptaciones de librerías (XGBoost $\ge$2, LightGBM $\ge$4, pandas~3) alteran detalles de
desempate y de parada temprana respecto de 2019; se priorizó preservar la \emph{lógica} de cada
pipeline sobre la igualdad bit a bit. Las diferencias con los puntajes históricos de Kaggle son
esperadas y se documentan en \cref{cap:resultados}; ninguna cambia la conclusión comparativa.

\chapter{Conclusiones y propuesta}
\label{cap:conclusiones}

\section{Conclusiones}

\begin{enumerate}
  \item La reproducción local de las mejores soluciones públicas de IEEE-CIS Fraud Detection es
        viable en hardware de consumidor (16 núcleos, RTX 2060 de 6\,GB) y entrega resultados
        cualitativamente idénticos a los de la competencia: el ordenamiento por AUC se conserva.
  \item El experimento de ablación demuestra que el rendimiento excepcional de la solución
        ganadora no proviene del algoritmo (idéntico en ambas variantes) ni del hardware, sino
        de la \textbf{ingeniería de variables orientada a entidades}: reconstruir el cliente
        implícito (\emph{magic UID}) y agregar su historial aporta $+0{,}0104$ de AUC local,
        más que cualquier cambio de modelo estudiado.
  \item La \textbf{estabilización temporal} de variables (normalización de columnas \texttt{D},
        filtrado por consistencia temporal) y la \textbf{validación temporal honesta}
        (\texttt{GroupKFold} mensual) son condiciones habilitantes: sin ellas, las ganancias no
        se transferirían al mes de prueba.
  \item El hardware acelera 4.4$\times$ el ciclo de experimentación sin mover el AUC: su valor es
        de productividad, no de precisión.
  \item Las técnicas de re-muestreo (SMOTE) presentes en el corpus analítico no forman parte de
        las soluciones ganadoras: con árboles y AUC, el desbalance se maneja vía objetivos
        probabilísticos y \texttt{scale\_pos\_weight}.
\end{enumerate}

\section{Propuesta de desarrollo}
\label{sec:futuro}

Sobre la evidencia recogida, la propuesta consiste en construir el pipeline siguiente, en orden
de retorno esperado:

\begin{enumerate}
  \item \textbf{Base}: XGBoost/LightGBM con las familias \texttt{C/D/V} seleccionadas por
        correlación y columnas \texttt{D} normalizadas (componentes ya verificados).
  \item \textbf{Entidades causales}: implementar el magic UID con agregaciones
        \emph{estrictamente históricas} (cada transacción usa solo estadísticas de
        transacciones anteriores a su \texttt{TransactionDT}), eliminando el carácter
        transductivo detectado en el protocolo original.
  \item \textbf{Ensemble}: combinar por promedio de rangos el XGB de entidades con un LightGBM
        equivalente ---el camino seguido por el equipo ganador (0{,}9459 en el tablero
        privado)--- y añadir el post-procesado de suavizado por UID.
  \item \textbf{Evaluación de negocio}: incorporar umbral de decisión y curva de utilidad
        (costo de bloqueo vs.\ pérdida por fraude) junto al AUC.
\end{enumerate}

Con ello, el proyecto transita de \emph{reproducir} las mejores técnicas públicas a
\emph{producir} una solución propia, superando la limitación metodológica principal detectada.


\appendix
\chapter{Reproducibilidad}
\label{cap:apendice}

\section{Estructura del proyecto}

\begin{verbatim}
PROYECTO1_PLANIFICA/
|-- data/raw/                 dataset oficial (CSV) + leaderboard publico
|-- codigos/
|   |-- ejecutados/           4 notebooks reproducidos (.ipynb + .py)
|   |-- descartados/          8 notebooks analizados sin ejecutar
|   `-- _RESUMEN_TECNICAS.md  inventario de tecnicas por notebook
|-- src/                      benchmarks reproducibles
|   |-- bench_01_inversion.py baseline XGB CPU
|   |-- bench_02_xhlulu.py    baseline XGB GPU
|   |-- bench_03_nroman.py    LightGBM + FE (TimeSeriesSplit)
|   |-- bench_04_cdeotte.py   XGB d12 con/sin magic UID
|   |-- download_notebooks.py descarga via API publica de Kaggle
|   |-- extract_notebooks.py  exportacion .ipynb -> .py
|   `-- make_ranking.py       consolidacion de resultados
|-- results/                  CSV y logs de cada benchmark + ranking
`-- informe_latex/            este informe (tectonic main.tex)
\end{verbatim}

\section{Entorno}

\begin{itemize}
  \item Linux; Python 3.12 en entorno \texttt{uv}; pandas 3.0, XGBoost 3.4,
        LightGBM 4.x, scikit-learn.
  \item CPU 16 núcleos; NVIDIA RTX 2060 6\,GB (driver 610).
  \item Datos: \texttt{kaggle competitions download -c ieee-fraud-detection}
        (requiere aceptar las reglas de la competencia).
\end{itemize}

\section{Ejecución de los experimentos}

\begin{verbatim}
uv venv --python 3.12 .venv
uv pip install pandas numpy scikit-learn lightgbm xgboost
uv pip install matplotlib
.venv/bin/python src/bench_01_inversion.py     # AUC 0.9191, 112 s
.venv/bin/python src/bench_02_xhlulu.py        # AUC 0.9227,  25 s (GPU)
.venv/bin/python src/bench_03_nroman.py        # AUC 0.9200, 462 s
.venv/bin/python src/bench_04_cdeotte.py       # 0.9375 -> 0.9479 (GPU)
.venv/bin/python informe_latex/scripts/generar_graficos.py
tectonic informe_latex/main.tex
\end{verbatim}

\section{Adaptaciones de versiones (2019 $\to$ 2026)}

\begin{itemize}
  \item \textbf{XGBoost $\ge$2}: \texttt{tree\_method='gpu\_hist'} se reemplaza por
        \texttt{tree\_method='hist', device='cuda'}; \texttt{early\_stopping\_rounds} y
        \texttt{eval\_metric} pasan al constructor del estimador.
  \item \textbf{LightGBM $\ge$4}: los argumentos \texttt{early\_stopping\_rounds} y
        \texttt{verbose\_eval} de \texttt{lgb.train} se sustituyen por los
        \emph{callbacks} \texttt{lgb.early\_stopping()} y \texttt{lgb.log\_evaluation()}.
  \item \textbf{pandas 3}: las cadenas leídas de CSV adoptan el dtype \texttt{str} (no
        \texttt{object}); las condiciones \texttt{dtype=='object'} deben reemplazarse por
        \texttt{pd.api.types.is\_numeric\_dtype(...)} negado. \texttt{LabelEncoder} se
        sustituye de forma segura por \texttt{pd.factorize(sort=True)} cuando hay tipos mixtos.
  \item \textbf{RAPIDS cuDF} (notebook GPU de Deotte): requiere $>$6\,GB de VRAM; no ejecutable
        en la RTX 2060, se analizó cualitativamente.
\end{itemize}

\section{Protocolo ético}

El estudio se limita a material público con licencia abierta (Apache 2.0 en los notebooks
citados). No se realizó ninguna submisión a la competencia ni uso de servicios externos con los
datos; toda la evidencia numérica proviene de ejecuciones locales.


\bibliographystyle{plain}
\bibliography{referencias}

\end{document}
